% --- Άσκηση 14730 (14730.tex) ---
\Askhsh{14730}{2}
\begin{ylh}{2}
    \item 2.1. Οι Πράξεις και οι Ιδιότητές τους
    \item 2.3. Απόλυτη Τιμή Πραγματικού Αριθμού
\end{ylh}
Δίνεται η παράσταση $A = |x - 2| + 3, x \in \mathbb{R}$.

\begin{alist}
\item Να βρείτε

\begin{rlist}
\item Tην τιμή της παράστασης A για $x = 2^3 - 3^2$. \monades{8}
\item Tις τιμές του x , ώστε να ισχύει $A = 5$. \monades{10}
\end{rlist}
\item Να εξετάσετε αν μπορεί η παράσταση $A$ να πάρει την τιμή 2. Να αιτιολογήσετε την απάντησή σας.

\monades{7}
\end{alist}
